% ---------------------------------------------------------------------------
% Laboratory 3 -- Process Scheduling
%
% Build:   make experiments     (produces results/*.csv)
%          pdflatex report.tex  (twice, for the table of contents)
% or simply: make report
%
% Every graph below reads directly from results/*.csv, so re-running the
% experiments on a different set of workload files regenerates the whole
% report without touching this source.
% ---------------------------------------------------------------------------
\documentclass[11pt,a4paper]{article}

\usepackage[margin=1in]{geometry}
\usepackage{booktabs}
\usepackage{graphicx}
\usepackage{listings}
\usepackage{xcolor}
\usepackage{amsmath}
\usepackage{pgfplots}
\usepackage[hidelinks]{hyperref}
\pgfplotsset{compat=1.17}

\pgfplotsset{
  bargraph/.style={
    ybar,
    width=0.95\linewidth, height=6.2cm,
    bar width=7pt,
    symbolic x coords={cpu-bound,io-bound,mixed,starve},
    xtick=data,
    enlarge x limits=0.18,
    ymajorgrids, grid style={dashed,gray!30},
    legend style={at={(0.5,1.03)}, anchor=south, legend columns=-1,
                  draw=none, font=\small, /tikz/every even column/.append style={column sep=8pt}},
  },
  linegraph/.style={
    width=0.95\linewidth, height=6.2cm,
    mark size=1.8pt,
    ymajorgrids, grid style={dashed,gray!30},
    legend style={at={(0.5,1.03)}, anchor=south, legend columns=-1,
                  draw=none, font=\small, /tikz/every even column/.append style={column sep=8pt}},
  },
}

\lstset{
  basicstyle=\ttfamily\footnotesize,
  columns=fullflexible,
  breaklines=true,
  frame=single,
  rulecolor=\color{gray!40},
  backgroundcolor=\color{gray!5},
  keywordstyle=\color{blue!60!black},
  commentstyle=\color{green!40!black}\itshape,
  language=C++,
}

\title{\textbf{Laboratory 3: Process Scheduling}\\[0.3em]
       \large A Simulation Study of FIFO, Round Robin and MLFQ}
\author{% TODO: replace with your name and roll number
        Name \\ Roll Number}
\date{\today}

\begin{document}
\maketitle

\begin{abstract}
\noindent
This report describes a discrete-event process scheduling simulator written in
C\texttt{++} and uses it to compare First-In-First-Out (FIFO), Round Robin (RR)
and the Multi-Level Feedback Queue (MLFQ) on four workloads with different
CPU/I/O characteristics, on both one and two processors. The headline result is
that no algorithm dominates: FIFO achieves the lowest average turnaround time on
three of the four workloads while delivering response times two orders of
magnitude worse than the alternatives, and it collapses completely on the fourth
workload, where a convoy of long jobs blocks short ones. MLFQ recovers almost all
of FIFO's turnaround performance while keeping response time comparable to RR
with the smallest quantum, and it does so without being told anything about job
lengths in advance. The periodic priority boost is shown to be a fairness
mechanism that trades response time for protection against starvation, rather
than an unambiguous improvement.
\end{abstract}

\tableofcontents
\newpage

% ===========================================================================
\section{Objective and Scope}

The assignment asks for a simulator that accepts a scheduling algorithm and a
workload description file, and reports turnaround time (average and maximum),
the run time of the simulator itself, and the resulting schedule. Part~I covers
FIFO, Round Robin with varying time quanta, and a three-level MLFQ with and
without a periodic priority boost. Part~II repeats all three on a
two-processor machine.

Alongside the two required metrics we also record \emph{response time} (first
time on a CPU minus arrival time), \emph{waiting time}, \emph{makespan},
\emph{CPU utilisation} and the number of \emph{dispatches}. These are not
required, but turnaround time alone cannot distinguish an algorithm that is
good from one that is merely finishing everything late at the same time, and
several of the conclusions below depend on that distinction.

% ===========================================================================
\section{Simulator Design}

\subsection{Representing a process}

A single structure holds both the static description read from the workload
file and the runtime state, as suggested in the assignment tip. The burst
sequence is stored exactly as it appears in the file, in one flat array: even
indices are CPU bursts and odd indices are I/O bursts. This removes the need to
keep two parallel arrays in step, and it makes ``advance to the next burst''
a single increment.

\begin{lstlisting}
struct Process {
    int  pid;  long long arrival;
    std::vector<long long> bursts;  // [0],[2],[4]... CPU ; [1],[3]... I/O

    std::size_t bi;        // burst currently being served
    long long   rem;       // time left in the current CPU burst
    int         cpuBurstNo;// 1-based ordinal, used when printing the schedule
    long long   first;     // first time on a CPU  -> response time
    long long   finish;    // completion time      -> turnaround time

    int       level;       // MLFQ queue index
    long long used;        // allotment consumed at the current level
    long long gen;         // boost generation (see section 2.4)
};
\end{lstlisting}

The schedule is accumulated as a list of contiguous slices per processor:

\begin{lstlisting}
struct Segment { int pid; int burstNo; long long start, end; }; // [start,end)
\end{lstlisting}

Slices are merged on insertion: if a process is preempted at the end of its
quantum and immediately redispatched onto the same CPU (which happens whenever
no other process is ready), the two slices are coalesced into one rather than
being reported as two adjacent entries.

\subsection{Queues and processor state}

Ready processes are held in \texttt{std::list<int>} queues --- one queue for
FIFO and RR, three for MLFQ. Waiting (blocked) processes are held in a
min-heap keyed by I/O completion time, so the next wake-up is always available
in $O(1)$ and insertion costs $O(\log n)$. Each processor has a currently
running process and a remaining-quantum counter, which is all the per-CPU state
the multiprocessor case needs.

\subsection{The simulation loop}

The simulator is event-driven rather than tick-by-tick. At each step it
computes the largest interval $\Delta$ during which no scheduling decision can
possibly change:
\[
\Delta = \min\Big(
  \min_{c \in \text{busy}} \text{rem}_c,\;
  \min_{c \in \text{busy}} \text{quantum-left}_c,\;
  t_{\text{next arrival}} - t,\;
  t_{\text{next I/O completion}} - t,\;
  t_{\text{next boost}} - t
\Big)
\]
and then advances every running process by $\Delta$ at once. The resulting
schedule is identical to what a per-tick loop would produce, but the cost is
proportional to the number of events rather than to the length of the
simulated timeline --- which matters, because simulator run time is one of the
metrics we are asked to report.

Within a single instant, events are applied in a fixed order so that the
simulation is deterministic and matches the textbook conventions:

\begin{enumerate}\itemsep2pt
  \item new arrivals are admitted;
  \item processes whose I/O has completed are returned to a ready queue;
  \item the periodic priority boost fires, if due;
  \item processes preempted at this same instant are re-queued \emph{behind}
        the two groups above (a process that has just used a full quantum
        should not jump ahead of one that has just become runnable);
  \item idle processors are filled from the highest-priority non-empty queue;
  \item under MLFQ only, a ready process of strictly higher priority displaces
        the lowest-priority running process.
\end{enumerate}

\subsection{Design decisions worth stating explicitly}

Several points are under-specified by the assignment and the textbook, so the
choices made here are recorded for reproducibility.

\paragraph{FIFO is non-preemptive.} A process runs until its CPU burst ends. On
returning from I/O it joins the back of the single ready queue, so FIFO here is
``first come, first served'' applied per CPU burst, not per process.

\paragraph{RR grants a fresh quantum on each dispatch.} A process preempted
after a full quantum goes to the tail of the queue and is given the full
quantum again next time.

\paragraph{MLFQ allotment is cumulative.} Following the revised rule in
\emph{Operating Systems: Three Easy Pieces}, a process is demoted once it has
consumed its allotment at a level \emph{in total}, regardless of how many times
it gave up the CPU in between. Without this rule a process can retain top
priority indefinitely by issuing an I/O request just before each quantum
expires. The allotment therefore survives both preemption and an I/O burst, and
is reset only on demotion or on a boost.

\paragraph{Higher priority preempts immediately.} When a process arrives or
returns from I/O into a queue above the one a running process belongs to, it
displaces that process at once rather than waiting for the current quantum to
drain. This is what makes MLFQ's response time so low in
Section~\ref{sec:resp}, and it is the direct implementation of MLFQ's first
rule.

\paragraph{The boost is $O(1)$.} A naive boost walks every process to reset its
level, costing $O(n)$ every 20 time units --- which is $O(n \cdot
\text{makespan})$ overall and, measurably, made MLFQ several times slower than
RR on large workloads. Instead the simulator increments a global
\emph{generation} counter and splices the lower queues onto the top queue
(constant time for a linked list); a process whose recorded generation is stale
is known to belong at level 0 with a fresh allotment, and pays that cost once,
lazily, the next time it is examined. The schedules produced are bit-identical
to the naive version, and MLFQ's run time at 1600 processes fell from
$55$\,ms to $6.7$\,ms.

\subsection{Output format}

The schedule is emitted exactly in the form of \texttt{schedule\_format.txt}:
a \texttt{CPU}$k$ header per processor, then one tab-separated line per slice
giving \texttt{P}\textit{pid}\texttt{,}\textit{burst} and the inclusive start
and end ticks. A burst of length $L$ beginning at tick $s$ is printed as
$s$ to $s+L-1$. The time origin is 0 by default; \texttt{--base 1} shifts it
if a 1-indexed timeline is preferred.

\begin{lstlisting}[language=]
$ ./scheduler fifo t1.txt
average turnaround time  : 13.0000
maximum turnaround time  : 14
simulator run time       : 0.0027 ms  (simulation only, no file I/O)

CPU0
P1,1	0	4
P2,1	5	7
P1,2	9	11
P2,2	12	14
\end{lstlisting}

\subsection{Verification}

The three algorithms were validated against schedules traced by hand on small
two-process workloads for which the correct answer can be written down
independently, including cases that exercise quantum expiry coinciding exactly
with burst completion, I/O completion coinciding with preemption, and MLFQ
demotion carried across an I/O burst. The $O(1)$ boost optimisation was
validated by checking that it reproduces the earlier implementation's metrics
to the last digit on every workload.

% ===========================================================================
\section{Experimental Setup}

Four workloads of 40 processes each were used. Should the course-supplied
workload files be used instead, every number and graph in this report
regenerates unchanged by re-running \texttt{make experiments}.

\begin{table}[h]
\centering
\begin{tabular}{llll}
\toprule
Workload & CPU bursts & I/O bursts & Character \\
\midrule
\texttt{cpu-bound} & 40--120 & 2--6   & long compute, negligible I/O \\
\texttt{io-bound}  & 2--8    & 10--40 & short compute, frequent blocking \\
\texttt{mixed}     & 5--60   & 3--25  & broad spread of both \\
\texttt{starve}    & \multicolumn{2}{l}{one hog (150--300) per three interactive jobs (1--4)}
                   & adversarial for FIFO \\
\bottomrule
\end{tabular}
\caption{Workload characteristics. All arrival times are nondecreasing.}
\end{table}

All runs are single-threaded on one machine; reported simulator run times are
the mean of five repetitions and exclude file parsing and schedule output, as
required. Note that all four workloads keep a single CPU essentially saturated
(utilisation $\geq 0.99$), so the comparisons below are all under heavy load.

% ===========================================================================
\section{Part I: Uniprocessor Results}

\subsection{Turnaround time}

\begin{figure}[h]
\centering
\begin{tikzpicture}
\begin{axis}[bargraph, ylabel={Average turnaround time},
             title={\small Average turnaround time, 1 processor}]
  \addplot table [x=workload,y=fifo]  {results/per_workload_1cpu.csv};
  \addplot table [x=workload,y=rr]    {results/per_workload_1cpu.csv};
  \addplot table [x=workload,y=mlfq]  {results/per_workload_1cpu.csv};
  \addplot table [x=workload,y=mlfqb] {results/per_workload_1cpu.csv};
  \legend{{FIFO},{RR ($q{=}2$)},{MLFQ},{MLFQ+boost}}
\end{axis}
\end{tikzpicture}
\caption{Average turnaround time per workload on one processor. Lower is
better.}
\label{fig:tat1}
\end{figure}

Figure~\ref{fig:tat1} contains the most important result in this report, and it
is not the expected one: \textbf{FIFO has the best average turnaround time on
three of the four workloads}. On \texttt{cpu-bound} it finishes at $4179$
against $5151$ for RR, an advantage of nearly 19\%.

This is not an anomaly, and it is worth spelling out because it contradicts the
intuition that time-sharing is simply better. Under saturation the total amount
of CPU work is fixed, so the makespan is essentially fixed too; what a
scheduler controls is the \emph{distribution} of completion times. FIFO runs
each job to completion and lets it leave, so early jobs finish very early.
Round robin interleaves everything, so every job's completion time drifts
toward the end of the run --- the average turnaround suffers even though no
individual job is treated unfairly. This is the same effect that makes
Shortest-Job-First optimal for average turnaround: finishing \emph{something}
early beats making uniform progress on everything.

The exception is \texttt{starve}, and it is decisive: FIFO records $3863$
against $1546$ for RR, a factor of 2.5 worse. This is the convoy effect. A few
processes with 150--300 unit CPU bursts arrive early, and every short
interactive job behind them waits for the whole convoy to drain. FIFO's
advantage on the other three workloads depends entirely on the accident that
short jobs happened not to queue behind long ones; when that assumption breaks,
FIFO has no mechanism to recover. \textbf{An algorithm that is 19\% better in
the good case and 150\% worse in the bad case is not a good trade} for a
general-purpose system, which is precisely the argument for MLFQ.

MLFQ tracks RR closely everywhere and matches it on \texttt{starve}
($1530$ with boost), without any advance knowledge of burst lengths --- it
infers them from observed behaviour.

\subsection{Response time}
\label{sec:resp}

\begin{figure}[h]
\centering
\begin{tikzpicture}
\begin{axis}[bargraph, ylabel={Average response time}, ymode=log,
             log origin=infty,
             title={\small Average response time, 1 processor (log scale)}]
  \addplot table [x=workload,y=fifo_r]  {results/per_workload_1cpu.csv};
  \addplot table [x=workload,y=rr_r]    {results/per_workload_1cpu.csv};
  \addplot table [x=workload,y=mlfq_r]  {results/per_workload_1cpu.csv};
  \addplot table [x=workload,y=mlfqb_r] {results/per_workload_1cpu.csv};
  \legend{{FIFO},{RR ($q{=}2$)},{MLFQ},{MLFQ+boost}}
\end{axis}
\end{tikzpicture}
\caption{Average response time per workload, one processor. Note the
logarithmic axis: the algorithms differ by more than three orders of
magnitude.}
\label{fig:resp1}
\end{figure}

Figure~\ref{fig:resp1} is the other half of the picture, and it reverses the
verdict. FIFO's average response time on \texttt{cpu-bound} is $1538.8$ time
units; MLFQ without boost achieves $0.3$. That is a factor of roughly $5000$.

MLFQ's response time is near-zero by construction: a newly arrived process
enters the top queue and immediately preempts whatever lower-priority process
is running, so it reaches a CPU within one time unit of arriving. RR with
$q{=}2$ must wait for the processes ahead of it in the queue to each consume a
quantum, giving $\approx 30$. FIFO must wait for entire CPU bursts to finish.

Taken together, Figures~\ref{fig:tat1} and~\ref{fig:resp1} say that FIFO buys
its turnaround advantage by making interactive work unusable, which is an
unacceptable trade on any machine a person is sitting in front of, and a
perfectly reasonable one on a batch cluster.

\subsection{Round robin: effect of the time quantum}

\begin{figure}[h]
\centering
\begin{tikzpicture}
\begin{axis}[linegraph, xlabel={Time quantum $q$}, ylabel={Turnaround time},
             xmode=log, log basis x=2, xtick={1,2,4,8,16,32,64},
             xticklabels={1,2,4,8,16,32,64},
             title={\small Round robin: turnaround time against quantum}]
  \addplot table [x=q,y=avg1] {results/rr_quantum.csv};
  \addplot table [x=q,y=avg2] {results/rr_quantum.csv};
  \addplot table [x=q,y=max1] {results/rr_quantum.csv};
  \addplot table [x=q,y=max2] {results/rr_quantum.csv};
  \legend{{avg (1 CPU)},{avg (2 CPU)},{max (1 CPU)},{max (2 CPU)}}
\end{axis}
\end{tikzpicture}
\caption{Mean over the four workloads. Average turnaround degrades with larger
quanta; maximum turnaround very slightly improves.}
\label{fig:rrq}
\end{figure}

\begin{figure}[h]
\centering
\begin{tikzpicture}
\begin{axis}[linegraph, xlabel={Time quantum $q$}, ylabel={Dispatches},
             xmode=log, log basis x=2, ymode=log,
             xtick={1,2,4,8,16,32,64}, xticklabels={1,2,4,8,16,32,64},
             title={\small Round robin: scheduling decisions against quantum}]
  \addplot table [x=q,y=sw1] {results/rr_quantum.csv};
  \legend{{dispatches (1 CPU)}}
\end{axis}
\end{tikzpicture}
\caption{Dispatch count falls roughly as $1/q$: $4247$ at $q{=}1$ down to
$215$ at $q{=}64$.}
\label{fig:rrsw}
\end{figure}

Average turnaround rises monotonically with the quantum, from $2676$ at
$q{=}1$ to $3136$ at $q{=}64$ --- a 17\% degradation. As $q$ grows, RR
approaches FIFO's behaviour of running a whole burst at a time, and inherits
FIFO's convoy problem on \texttt{starve}; the small quantum wins because it
approximates SJF, letting short jobs slip out between slices of long ones.

Maximum turnaround moves in the opposite direction, but barely ($4224
\rightarrow 4188$, under 1\%). The last process to finish is limited by total
system work, which no quantum choice changes.

The important caveat is that this simulator, like the assignment's model,
charges \textbf{nothing for a context switch}. Figure~\ref{fig:rrsw} shows
$q{=}1$ requires $4247$ dispatches against $215$ at $q{=}64$ --- twenty times
as many. On real hardware, at a plausible $1$--$10\,\mu$s per switch including
cache and TLB damage, that difference is exactly what makes $q{=}1$
impractical and pushes real systems to quanta in the millisecond range. The
graph showing small quanta winning should therefore be read as ``small quanta
win \emph{on the modelled dimensions}''; $q{=}2$ or $q{=}4$ is the defensible
choice, since it captures nearly all the turnaround benefit at a quarter to an
eighth of the switching cost.

\subsection{MLFQ: effect of the periodic priority boost}

\begin{table}[h]
\centering
\begin{tabular}{lrrrr}
\toprule
& \multicolumn{2}{c}{Avg.\ turnaround} & \multicolumn{2}{c}{Avg.\ response} \\
\cmidrule(lr){2-3}\cmidrule(lr){4-5}
Workload & no boost & boost 20 & no boost & boost 20 \\
\midrule
\texttt{cpu-bound} & 5162.9 & 5155.8 & \phantom{0}0.3 & 29.1 \\
\texttt{io-bound}  & \phantom{0}992.6 & \phantom{0}991.3 & \phantom{0}0.8 & 25.0 \\
\texttt{mixed}     & 3032.1 & 3042.6 & \phantom{0}1.1 & 29.4 \\
\texttt{starve}    & 1564.4 & \textbf{1529.6} & \phantom{0}5.4 & 24.6 \\
\bottomrule
\end{tabular}
\caption{Boost every 20 time units versus no boost, one processor.}
\label{tab:boost}
\end{table}

The boost's effect on turnaround time is small and not uniformly positive: it
helps on \texttt{starve} (2.2\% better) and on the two homogeneous workloads by
a fraction of a percent, and it is very slightly \emph{worse} on
\texttt{mixed}. Anyone expecting the boost to be a clear win will not find one
here.

Its effect on response time is large and negative: from $0.3$--$5.4$ without
boost to $24.6$--$29.4$ with it, roughly an order of magnitude. The mechanism
is direct. Without a boost, long-running processes sink to the bottom queue and
stay there, so the top queue is occupied only by genuinely new or interactive
work and a new arrival is served instantly. Every boost re-promotes the CPU
hogs to the top queue, where they compete with new arrivals on equal terms
until they are demoted again.

The boost is therefore best understood not as an optimisation but as
\textbf{insurance}. The quantity it protects is not visible in either average
above --- it is the worst case for an individual long-running process. Under
these saturated workloads, demoted processes still make progress whenever the
upper queues empty, so no process starves outright and the insurance never has
to pay out; its premium (worse response time) is all that shows up. On a
workload with a permanent stream of short interactive jobs the bottom queue
would never be reached and a no-boost MLFQ would starve long jobs indefinitely.
A period of 20 is also aggressive relative to these bursts; a longer period
would collect a smaller premium for nearly the same protection.

% ===========================================================================
\section{Part II: Two Processors}

\begin{figure}[h]
\centering
\begin{tikzpicture}
\begin{axis}[bargraph, ylabel={Average turnaround time},
             title={\small Average turnaround time, 2 processors}]
  \addplot table [x=workload,y=fifo]  {results/per_workload_2cpu.csv};
  \addplot table [x=workload,y=rr]    {results/per_workload_2cpu.csv};
  \addplot table [x=workload,y=mlfq]  {results/per_workload_2cpu.csv};
  \addplot table [x=workload,y=mlfqb] {results/per_workload_2cpu.csv};
  \legend{{FIFO},{RR ($q{=}2$)},{MLFQ},{MLFQ+boost}}
\end{axis}
\end{tikzpicture}
\caption{Two processors. The ranking of the algorithms is unchanged from
Figure~\ref{fig:tat1}.}
\label{fig:tat2}
\end{figure}

\begin{figure}[h]
\centering
\begin{tikzpicture}
\begin{axis}[bargraph, ylabel={Speedup (makespan$_1$/makespan$_2$)},
             ymin=1.5, ymax=2.1,
             title={\small Speedup from adding a second processor}]
  \addplot table [x=workload,y=fifo] {results/speedup.csv};
  \addplot table [x=workload,y=rr]   {results/speedup.csv};
  \addplot table [x=workload,y=mlfq] {results/speedup.csv};
  \legend{{FIFO},{RR ($q{=}2$)},{MLFQ+boost}}
\end{axis}
\end{tikzpicture}
\caption{Speedup is near-linear everywhere except the I/O-bound workload.}
\label{fig:speedup}
\end{figure}

Two processors roughly halve every turnaround figure and leave the relative
ranking of the algorithms untouched, which is itself worth noting: the
conclusions of Part~I are not artefacts of a single CPU.

Speedup (Figure~\ref{fig:speedup}) is between $1.98$ and $2.00$ on
\texttt{cpu-bound}, \texttt{mixed} and \texttt{starve} --- essentially linear.
These workloads keep both processors saturated (utilisation $0.99$), so
doubling the supply of CPU halves the time to drain a fixed amount of work.

\texttt{io-bound} is the exception, at $1.86$--$1.90$. Its utilisation falls
from $0.99$ on one processor to $0.89$--$0.94$ on two: with short CPU bursts
and long I/O bursts, there are simply not always two runnable processes
available, and a processor sits idle. The bottleneck has shifted from the CPU
to the I/O device, and adding processors cannot help with that. This is the
practical form of Amdahl's argument --- the achievable speedup is governed by
the fraction of the workload that genuinely needs a CPU.

Two implementation notes. First, the ready queues are global and shared between
processors, so a process may migrate freely between CPU 0 and CPU 1 between
slices; a real scheduler would weigh cache affinity against load balance, which
this model does not capture and which would erode the near-linear speedup
above. Second, the model charges nothing for contention on the shared ready
queue, which on a real multiprocessor is a lock that becomes a scalability
bottleneck well before the processor count gets large --- the reason production
schedulers use per-CPU run-queues with periodic work stealing.

% ===========================================================================
\section{Simulator Run Time}

\begin{figure}[h]
\centering
\begin{tikzpicture}
\begin{axis}[linegraph, xlabel={Number of processes}, ylabel={Run time (ms)},
             xmode=log, ymode=log, log basis x=2,
             title={\small Simulator run time, excluding file I/O}]
  \addplot table [x=n,y=fifo] {results/scale.csv};
  \addplot table [x=n,y=rr]   {results/scale.csv};
  \addplot table [x=n,y=mlfq] {results/scale.csv};
  \legend{{FIFO},{RR ($q{=}2$)},{MLFQ+boost}}
\end{axis}
\end{tikzpicture}
\caption{Run time against workload size on log--log axes; all three curves are
straight with unit slope, i.e.\ linear in the number of events.}
\label{fig:scale}
\end{figure}

All three algorithms scale linearly in the number of simulated events. FIFO is
roughly ten times cheaper than the other two at every size --- $0.65$\,ms
against $6.7$\,ms at 1600 processes --- for a structural reason: FIFO generates
one event per CPU burst, whereas RR and MLFQ generate one per quantum, so a
burst of length $L$ costs $\lceil L/q \rceil$ events instead of one. The
measured cost is therefore a direct proxy for the number of scheduling
decisions the algorithm makes, and it reproduces Figure~\ref{fig:rrsw} from a
completely different direction.

This is also where the boost optimisation of Section~2.4 shows up. With the
naive $O(n)$ boost, MLFQ's curve bent upward and reached $55$\,ms at 1600
processes; it now lies on top of RR's. Had that not been fixed, the graph would
have appeared to show that MLFQ is intrinsically expensive to schedule, which
is a statement about one implementation rather than about the algorithm.

% ===========================================================================
\section{Conclusions}

\begin{enumerate}\itemsep3pt
  \item \textbf{No algorithm wins on both metrics.} FIFO gives the best average
        turnaround on three of four workloads and the worst response time by a
        factor of thousands. Any single-number comparison of schedulers should
        be treated with suspicion.

  \item \textbf{FIFO's advantage is fragile.} It depends on short jobs not
        queueing behind long ones. On the workload where they do, FIFO is 2.5
        times worse than everything else. Predictable mediocrity beats
        unpredictable excellence in a general-purpose scheduler.

  \item \textbf{Smaller quanta improve turnaround in this model, but the model
        is generous.} Average turnaround degrades 17\% from $q{=}1$ to
        $q{=}64$, while dispatches fall twentyfold. With any realistic context
        switch cost the optimum moves sharply upward; $q{=}2$--$4$ is the
        honest recommendation here.

  \item \textbf{MLFQ delivers most of both.} It approaches FIFO's turnaround
        and beats RR's response time, without advance knowledge of job lengths,
        by inferring behaviour from observed CPU usage.

  \item \textbf{The priority boost is insurance, not an optimisation.} It costs
        an order of magnitude in response time and returns at most 2\% in
        turnaround on these workloads. Its value lies in a worst case that
        these saturated workloads never trigger.

  \item \textbf{A second processor gives near-linear speedup only while the CPU
        is the bottleneck.} Speedup is $\approx 2.0$ on CPU-heavy workloads and
        $1.86$ on the I/O-bound one, where utilisation drops to $0.89$.
\end{enumerate}

\paragraph{Limitations.} Context switches, cache affinity, migration cost and
ready-queue contention are all free in this model; I/O devices are unbounded and
never queue; and the workloads are synthetic. Each of these would move the
results toward larger quanta, fewer preemptions and sublinear multiprocessor
speedup. The qualitative conclusions --- the convoy effect, the
turnaround/response tension, MLFQ's inference of job type --- do not depend on
them.

% ===========================================================================
\appendix
\section{Reproducing These Results}

\begin{lstlisting}[language=bash]
make                 # build ./scheduler
make gen-workloads   # generate the four synthetic workloads
make experiments     # 80+ simulator runs -> results/*.csv
make report          # pdflatex x2 -> report.pdf

# individual runs
./scheduler fifo workloads/mixed.txt
./scheduler rr   workloads/mixed.txt --quantum 4 --cpus 2
./scheduler mlfq workloads/mixed.txt --quantum 2 --boost 20 --schedule sched.txt
\end{lstlisting}

\begin{table}[h]
\centering
\begin{tabular}{ll}
\toprule
Option & Meaning \\
\midrule
\texttt{--quantum <q>}  & time quantum for RR / MLFQ (default 2) \\
\texttt{--cpus <n>}     & number of processors (default 1) \\
\texttt{--boost <b>}    & MLFQ boost period; 0 disables it (default 0) \\
\texttt{--levels <l>}   & number of MLFQ queues (default 3) \\
\texttt{--base <b>}     & time origin used when printing the schedule \\
\texttt{--schedule <f>} & write the schedule to file \texttt{f} \\
\texttt{--per-process}  & print a per-process turnaround/response table \\
\texttt{--csv}          & emit one CSV record instead of a human report \\
\bottomrule
\end{tabular}
\caption{Command-line options.}
\end{table}

\end{document}
